% Complete local realizer constructions and separation proofs.

\begin{annexitem}{annex:p04-wc-basic}{thm:p04-wc-basic}{Composition, representation translations and the lost-input example}
\begin{proof}
For reflexivity take \(K=\mathrm{id}\), with the postprocessor projecting
the oracle output; for strong reflexivity use both identity operators.
Suppose \(e\leq_{\mathrm W} f\) has preprocess/postprocess operators
\(K_1,H_1\), and \(f\leq_{\mathrm W} g\) has \(K_2,H_2\).
The combined preprocessor is \(K_2K_1\). On original input \(p\)
and final oracle answer \(q\), the combined postprocessor computes
\[
 H_1\bigl(\langle p,H_2(\langle K_1(p),q\rangle)\rangle\bigr).
\]
For every realizer of \(g\), the inner computation is an allowed
\(f\)-answer on \(K_1(p)\), and the outer one an allowed \(e\)-answer.
All promised domains are preserved, and the operators are computable
by buffered stream composition. The strong case simply composes
postprocessors \(H_1H_2\) and preprocessors \(K_2K_1\).
Discarding the original input in an ordinary postprocessor proves
strong-to-ordinary inclusion. Substituting an existing computable or
continuous oracle realizer proves preservation by the same compositions.

For representation invariance, let \(a_X\) translate new \(X\)-names
to old ones, \(b_Y\) old \(Y\)-names to new ones,
\(b_U\) old \(g\)-input names to new ones, and \(a_V\) new
\(g\)-output names to old ones. From old reduction operators \(H,K\),
use preprocessor \(b_UKa_X\) and ordinary postprocessor
\[
 \langle p,q\rangle\longmapsto
 b_YH\langle a_X(p),a_V(q)\rangle.
\]
Every new realizer thereby induces an old valid oracle answer, so the
decoded membership in \(f(x)\) is preserved. In the strong case use
postprocessor \(b_YHa_V\). The reverse translators prove the converse;
no continuous-only equivalence is asserted to preserve effectivity.

The ordinary identity-to-constant reduction submits a fixed computable
oracle input and returns its original \(p\), ignoring the answer.
For a hypothetical strong reduction, select the valid constant realizer
that outputs the same canonical zero name on every input. The strong
postprocessor would then produce the same Baire stream on every \(p\),
and hence cannot realize the identity at two distinct streams.
The quantifier over every realizer is essential to this contradiction.
\end{proof}
\end{annexitem}

\begin{annexitem}{annex:p04-wc-cylinder}{thm:p04-wc-cylinder}{The input retained in a cylindrical oracle}
\begin{proof}
Suppose \(f\leq_{\mathrm W} g\) with operators \(H,K\). For input
\(\langle r,p\rangle\) to \(\mathrm{id}\times f\), submit
\(\langle\langle r,p\rangle,K(p)\rangle\) to
\(\mathrm{id}\times g\). Every valid oracle answer has first
component exactly \(\langle r,p\rangle\), because the identity's
representation is identity. Its second component is some valid
\(g\)-answer name \(q\). The strong postprocessor returns
\(\langle r,H\langle p,q\rangle\rangle\).
This works for every allowed oracle answer and proves the forward direction.

Conversely, let \(K',H'\) strongly reduce the two cylindrifications.
Given \(p\), form \(\langle p,p\rangle\) and split
\(K'\langle p,p\rangle=\langle a(p),b(p)\rangle\).
Use \(b(p)\) as the ordinary oracle input to \(g\).
Given its answer \(q\), form \(\langle a(p),q\rangle\), apply
\(H'\), and project the \(f\)-answer component. Taking the
cylindrical oracle to be \(\mathrm{id}\times G\) proves correctness
for every \(G\vdash g\).

One always has \(f\leq_{\mathrm{sW}}\mathrm{id}\times f\): supply a fixed
first identity component and project the second output. If \(g\) is
a cylinder, the forward direction and this observation give
\(f\leq_{\mathrm{sW}}\mathrm{id}\times f
\leq_{\mathrm{sW}}\mathrm{id}\times g\leq_{\mathrm{sW}}g\).
The converse is strong-to-ordinary inclusion.
\end{proof}
\end{annexitem}

\begin{annexitem}{annex:p04-wc-parallel}{thm:p04-wc-parallel}{All coordinates, flattening and interleaving}
\begin{proof}
For \(f\leq_{\mathrm{sW}}\widehat f\), duplicate the one supplied name
\(p\) into every input coordinate and project the zeroth answer.
No externally chosen instance is needed. If the promise is empty,
there is no promised input to check.

For monotonicity, preprocess every component \(p_i\) by the given
\(K\). A realizer of \(\widehat g\) returns a tuple \((q_i)\)
with each \(q_i\) an allowed answer on \(K(p_i)\).
Apply \(H\langle p_i,q_i\rangle\) at every coordinate, or
\(H(q_i)\) in the strong case. The reduction requirement for every
realizer guarantees correctness for each allowed answer, even if the
parallel oracle correlates its choices between coordinates. This can
also be seen by extending any one allowed answer at the one input name
to an otherwise arbitrary realizer. Each requested finite tuple-entry
is computed by a finite component computation, giving a productive
countable stream output.

For double parallelization flatten the matrix of instance names using
the fixed pairing \(\langle i,j\rangle\), and rearrange the resulting
answer stream into the original matrix. Conversely partition a single
list into the paired matrix, call the double problem and flatten its
answers. These inverse computable permutations work in both directions
without accessing the original input in postprocessing.
For the product identity, zip lists of \(f\)- and \(g\)-instances
into a list of pairs and unzip the answer pairs; the inverse does the
opposite. For two copies of \(\widehat f\), interleave their input
lists at even and odd coordinates and split their answer list accordingly.
In the reverse direction duplicate a list and retain one result. These
constructions establish every displayed strong equivalence.
\end{proof}
\end{annexitem}

\begin{annexitem}{annex:p04-wc-finite}{thm:p04-wc-finite}{Length tags and the zero-length case}
\begin{proof}
One-element lists give extensiveness. For monotonicity preprocess each
of the finitely many supplied instances with \(K\); retain the length
tag. On each oracle-answer component use the original component name
and \(H\) as in the countable proof. When the tag is zero, no oracle
instance is requested and the answer is the empty list.

For \((f^*)^*\leq_{\mathrm W} f^*\), read the finitely many sublist lengths,
concatenate all their instances, and submit that one finite list.
The ordinary postprocessor reads the retained original length tags and
partitions the answer list into the original sublists. If the total
length is zero, it still recovers the correct number of empty sublists
from those tags. The reverse reduction places a list inside a one-list
outer container and projects the first returned sublist.

If a computable promised instance \(x_0\) is supplied as a fixed
constant, extend any finite list by indefinitely many copies of its name
to give a countable list. Call \(\widehat f\) and retain the first
originally tagged number of answers. This includes empty lists and
proves the last reduction. No such padding premise has been assumed
for a problem without a computable point in its domain. The empty-list
constant is ordinarily equivalent to the identity by ignoring the
oracle and computing the required output; it is not strongly equivalent
to the identity by the constant-realizer obstruction already proved.
\end{proof}
\end{annexitem}

\begin{annexitem}{annex:p04-wc-finite-choice}{thm:p04-wc-finite-choice}{Both finite translations and the source answer-label shift}
\begin{proof}
Given a negative name of a nonempty subset of \(\{0,\ldots,n\}\),
for each \(i\in\{1,\ldots,n+1\}\) issue a stream \(p_i\)
that remains zero until point \(i-1\) is excluded, and contains
a nonzero marker when that exclusion is observed. It is the zero
stream exactly when the point remains in the set. At least one
component is therefore zero, satisfying the MLPO promise. Its
answer \(i\) yields the point \(i-1\) without the original input.
Conversely, from an MLPO tuple enumerate point \(i-1\) into the
complement as soon as a nonzero coordinate of \(p_i\) is observed.
The remaining finite set is nonempty by the promise. Any chosen
point \(j\) gives the allowed tuple answer \(j+1\).
All finite tests and enumeration scheduling are computable.
For \(n=1\) these translations, composed with the separate
binary-choice/single-marker translations below, prove strong
equivalence of the two source LLPO presentations. Their different
domains and answer labels have not been identified literally.
\end{proof}
\end{annexitem}

\begin{annexitem}{annex:p04-wc-binary-choice}{thm:p04-wc-binary-choice}{Single markers, discrete choice and the LPO variant bridge}
\begin{proof}
For an LLPO input \(p\), negatively name the subset of \(\{0,1\}\)
by excluding \(i\) if a nonzero coordinate \(p(2n+i)\) is found.
Enumeration can repeat discoveries and use padding; at most one of
the two points is excluded. The remaining set is exactly
\(\mathsf{LLPO}(p)\), so copying a choice answer gives a strong reduction.

Conversely, a nonempty negative binary set can exclude at most one bit
\(i\). Scan its name in stages. If the first exclusion occurs, schedule
one marker at a future stream position of parity \(i\), then output
zeros forever; before the event output zeros. Repeated exclusions are
ignored. The resulting LLPO instance has at most one nonzero entry,
and its permitted answers are exactly the unexcluded bits. Its answer
can again be copied without the original input.

To reduce LLPO to the native LPO, form
\(q(n)=1\) when \(p(2n)=0\), and \(q(n)=0\) otherwise.
If \(\mathsf{LPO}(q)=1\), all even entries of \(p\) vanish, so
output zero. If its value is zero, an even nonzero entry occurred;
the LLPO promise makes every odd entry zero, so output one.
The postprocessor is just the complement of the returned bit.

For \(\mathsf{LPO}_0\leq_{\mathrm{sW}}\mathsf{LPO}\), use
\(q(n)=1\) if \(p(n)=0\), and zero otherwise: its native LPO
answer is one exactly when \(p\) is the zero stream.
For the converse use the same coordinatewise transform; its zero-stream
test answer is one exactly when the original \(p\) has no zero entry.
These are computable transforms on discrete stream entries, not
decisions about an entire stream at once. Natural-output name tails
are ignored and a canonical name of the returned bit is produced.
\end{proof}
\end{annexitem}

\begin{annexitem}{annex:p04-wc-omniscience-obstruction}{thm:p04-wc-omniscience-obstruction}{Finite prefixes cannot settle the promised discrete output}
\begin{proof}
At the input \(1^{\mathbb N}\), a continuous LPO realizer must
issue a natural-output first entry equal to one. It fixes that entry
on some finite prefix. Extend that prefix by a zero and then ones;
the native LPO answer must now be zero, a contradiction.
For LLPO consider the input \(0^{\mathbb N}\). Whatever allowed
first output bit \(i\) the realizer chooses, continuity fixes it on
a finite zero prefix. Extend that prefix by a single nonzero entry
at a later position of parity \(i\), keeping all other entries zero.
This is still a promised input, but its sole allowed answer is
\(1-i\), again a contradiction. This addresses all possible
name-dependent choices at the initial input. Finite-use continuity
rules out computable realizers, and the binary-choice equivalence
transfers the obstruction.
\end{proof}
\end{annexitem}

\begin{annexitem}{annex:p04-wc-natural-limits}{thm:p04-wc-natural-limits}{The scalar encoding and the whole-stream counterboundary}
\begin{proof}
Given a negative name of a nonempty \(A\subseteq\mathbb N\),
let \(a_s\) be the least number not excluded among the first \(s\)
name entries. This finite search always terminates: a finite list
excludes only finitely many numbers. Let \(a\) be the least member
of \(A\). It is never excluded, and each of the finitely many smaller
numbers is eventually excluded. Hence \(a_s=a\) for all sufficiently
large \(s\). Calling \(\lim_{\mathbb N}\) and returning its answer
therefore gives a strong reduction \(\mathsf{C}_{\mathbb N}\leq_{\mathrm{sW}}\lim_{\mathbb N}\).

For an eventually constant numerical sequence \((a_s)\), negatively
name the following subset, using the fixed numerical pairing:
\[
 B=\{\langle s,v\rangle:(\forall t\geq s)\ a_t=v\}.
\]
The complement is enumerable: for every \((s,v)\) search for a
finite witness \(t\geq s\) with \(a_t\ne v\), dovetailing the
searches and padding when needed. The set is nonempty by eventual
constancy. Every choice answer decodes a pair with \(v\) equal to
the final value, so projecting \(v\) proves the reverse strong reduction.

For a sequence of entire streams \((p_s)\) promised eventually
constant, enumerate the complement of
\[
 S=\{s:(\forall t\geq s)(\forall j)\ p_t(j)=p_s(j)\}.
\]
A violation has a finite coordinate witness, and the dovetailed
enumeration is computable. Eventual whole-stream constancy makes
\(S\) nonempty. A natural-choice answer \(s\) permits the ordinary
postprocessor, which retains the original input, to return \(p_s\).
Conversely, the least-unexcluded approximations \(a_s\) above
give eventually constant entire streams \((a_s,0,0,\ldots)\).
Their \(\lim_\Delta\) output first entry is in \(A\), proving
\(\mathsf{C}_{\mathbb N}\leq_{\mathrm{sW}}\lim_\Delta\) and ordinary equivalence.

For failure of the reverse strong reduction choose a noncomputable
Baire stream \(p\), which exists by diagonalization over finite
programs. Submit the constant sequence \(p,p,p,\ldots\) to any
purported strong reduction of \(\lim_\Delta\) to \(\mathsf{C}_{\mathbb N}\).
Its preprocessor must name a nonempty natural set. Choose any allowed
number \(n\), and a valid oracle realizer returning its canonical
computable name at that particular input. The computable strong
postprocessor on that fixed computable name would return the required
stream \(p\), making \(p\) computable, a contradiction. The
ordinary reduction avoids this failure precisely by retaining the input.
\end{proof}
\end{annexitem}

\begin{annexitem}{annex:p04-wc-parallel-limit}{thm:p04-wc-parallel-limit}{All stabilization questions and both parallelized equivalences}
\begin{proof}
For a list of native LPO instances \((p_i)\), produce Baire streams
\[
 u_s(i)=\begin{cases}0,&(\exists j\leq s)\ p_i(j)=0,\\1,&\text{otherwise}.
 \end{cases}
\]
Every coordinate changes at most once and converges to
\(\mathsf{LPO}(p_i)\). The limit output is exactly the answer list;
tupling or padding converts between its Baire stream and the sequence
of natural answer names. This is a strong reduction.

Conversely, for coordinatewise convergent \((p_s)\), ask native LPO
in parallel for every triple \((n,k,a)\) on the stream
\[
 q_{n,k,a}(t)=
 \begin{cases}1,&p_{k+t}(n)=a,\\0,&p_{k+t}(n)\ne a.
 \end{cases}
\]
Its answer is one exactly when all values from stage \(k\) onward
equal \(a\). For each \(n\), eventual constancy ensures that at
least one pair \((k,a)\) has answer one. The postprocessor searches
the returned answer table until it finds one, then outputs \(a\)
at coordinate \(n\). Every successful pair gives the same eventual
value; the search terminates and uses no original \(p_s\).
This proves \(\lim\leq_{\mathrm{sW}}\widehat{\mathsf{LPO}}\).

The strong equivalence \(\mathsf{C}_{\mathbb N}\equiv_{\mathrm{sW}}\lim_{\mathbb N}\)
parallelizes by the closure theorem. A list of numerical sequences
\((a_{i,s})_s\), one per \(i\), is computably transposed into a
sequence of Baire streams \(u_s(i)=a_{i,s}\). Their coordinatewise
limit is the list of the scalar limits. In the reverse direction take
one numerical sequence per coordinate of a Baire-limit instance.
Both transpositions are inverse computable name rearrangements with
the identical eventual-constancy promises. This proves the second
strong equivalence without substituting whole-stream convergence.
\end{proof}
\end{annexitem}

\begin{annexitem}{annex:p04-wc-cantor-tree}{thm:p04-wc-cantor-tree}{A characteristic tree from an exterior enumeration}
\begin{proof}
First recall binary-tree compactness. If a prefix-closed binary tree
is infinite, it has nodes of arbitrarily large depth, since each
bounded set of levels is finite. Starting at the empty word, choose
at each step a child with extensions of arbitrarily large depth.
One such child exists because there are only two children. This
classical inductive choice gives an infinite path. Equivalently,
an open-cylinder cover of Cantor space with no finite subcover would
permit successive child cylinders with no finite subcover; their nested
prefixes give a point missed by the cover, a contradiction. Thus
Cantor space and each of its closed subsets are compact. These
existence arguments are not claimed as effective path algorithms.

Let \((v_j)\) be the negative name entries, each either a padded
empty exclusion or a finite word to exclude. Define
\[
 T=\{s\in\{0,1\}^{<\mathbb N}:
       \text{no excluded word among the first }|s|
       \text{ entries is a prefix of }s\}.
\]
Membership uses only finitely many entries and finite-word comparisons,
so it gives a characteristic-tree name uniformly. It is prefix-closed:
a forbidden prefix of a shorter word among its tested entries would
also forbid any extension. Every \(b\in A\) has all prefixes in
\(T\). Conversely, if \(b\notin A\), some excluded word \(v_j\)
is a prefix of \(b\); a sufficiently long prefix of \(b\), longer
than both \(j\) and \(|v_j|\), fails the test. Hence \([T]=A\).
Nonempty \(A\) makes \(T\) infinite, so WKL is a valid request,
and its binary answer is already an answer to Cantor choice.

Conversely, from a characteristic-tree name enumerate every finite
word \(s\notin T\), excluding its cylinder. A binary stream is
outside \([T]\) exactly when it has one of these prefixes.
Thus the enumerated closed set is exactly \([T]\), nonempty for an
infinite binary tree. An answer is already a path. Both postprocessors
are identity on binary streams, proving strong equivalence.
\end{proof}
\end{annexitem}

\begin{annexitem}{annex:p04-wc-cantor-llpo}{thm:p04-wc-cantor-llpo}{Coordinate products and guarded child queries}
\begin{proof}
Given promised LLPO instances \((p_i)\), consider
\[
 A=\{b\in\mathcal C:(\forall i)(\forall n)\ p_i(2n+b(i))=0\}.
\]
Each coordinate permits at least one bit, so \(A\ne\varnothing\).
For example, classically choose at coordinate \(i\) the least
permitted bit; this establishes nonemptiness, not its computability.
The complement is an enumerable union of cylinders: enumerate a
finite word \(s\) when some \(i<|s|\) and \(n\) witness
\(p_i(2n+s(i))\ne0\). Every point outside \(A\) has a finite
witness and such a prefix, and no point of \(A\) is excluded.
An answer \(b\in A\) supplies a valid LLPO answer \(b(i)\) in
each coordinate; padding the scalar names is computable and uses
only \(b\).

For the converse, given negative nonempty \(A\), the assertion
\(A\cap[s]=\varnothing\) is semidecidable. It says that the
excluded cylinders cover the compact cylinder \([s]\), and a finite
subcover must suffice. Given finitely many words, coverage of \([s]\)
is decidable: take a common depth at least all their lengths and
\(|s|\), and check every binary extension of \(s\) at that depth
for an excluded prefix. Search the successive finite negative-name
prefixes for this finite test.

For every finite word \(s\), run the two emptiness tests for
\(A\cap[s0]\) and \(A\cap[s1]\) in finite stages. At the first
stage with a verified empty child, choose one verified bit \(i\),
schedule a single nonzero marker at a later position of parity \(i\),
and ignore all further discoveries. Until then and afterwards issue
zeros. This computably supplies one LLPO input \(q_s\) for every
node; all have at most one nonzero entry, including nodes for which
both children are empty. All rows are generated independently in
parallel by finite computations from the negative name.

If \(A\cap[s]\ne\varnothing\), both children cannot be empty.
If one is empty, its test eventually succeeds and the LLPO answers
exclude exactly that child; if neither is empty, either answer is
valid. Thus every allowed answer \(i\) at such a node satisfies
\(A\cap[si]\ne\varnothing\).
From the parallel answer table start at the empty word, read its
answer, append that bit, and repeat at the new prefix. By induction
every resulting prefix-cylinder meets \(A\). The resulting binary
stream belongs to \(A\): if it belonged to an excluded cylinder,
a sufficiently long prefix would be contained in that cylinder and
could not meet \(A\). The postprocessor reads only the answer table,
not the original closed-set name. This proves the reverse strong reduction.
\end{proof}
\end{annexitem}

\begin{annexitem}{annex:p04-wc-omniscience-cylinders}{thm:p04-wc-omniscience-cylinders}{Single-marker coding of the retained input}
\begin{proof}
For each Baire input \(p\) and pair \((k,m)\), submit the constant
native LPO stream equal to zero if \(p(k)=m\) and equal to one
otherwise. The answer table is zero exactly at \(m=p(k)\).
Searching its \(k\)-th row for a zero recovers \(p(k)\).
The searches terminate for every coordinate, proving
\(\mathrm{id}\leq_{\mathrm{sW}}\widehat{\mathsf{LPO}}\).

For LLPO instead submit a row with a single marker at position one
when \(p(k)=m\), or a single marker at position zero otherwise.
All other entries are zero. These rows satisfy the at-most-one-nonzero
promise, and their unique answers are zero precisely when \(p(k)=m\).
The same search recovers \(p\) from every parallel answer, proving
\(\mathrm{id}\leq_{\mathrm{sW}}\widehat{\mathsf{LLPO}}\).

For either parallel problem \(Q\), interleave the input-recovery
queries just described with the supplied instance list for \(Q\).
A single \(Q\) answer can be split into both tables. The first
recovers the original independent identity input; the second is a
valid answer to the supplied \(Q\)-instance. Thus
\(\mathrm{id}\times Q\leq_{\mathrm{sW}}Q\).
The reverse reduction supplies any fixed identity first component
and projects the second answer, so \(Q\) is a cylinder.
Strong equivalence preserves this property: if \(Q\equiv_{\mathrm{sW}}R\),
products of the reduction operators give
\(\mathrm{id}\times R\leq_{\mathrm{sW}}\mathrm{id}\times Q
\leq_{\mathrm{sW}}Q\leq_{\mathrm{sW}}R\).
Applying the established strong equivalences proves the other listed
cylinders and their ordinary/strong coincidence.
\end{proof}
\end{annexitem}

\begin{annexitem}{annex:p04-wc-real-llpo}{thm:p04-wc-real-llpo}{Decaying real magnitudes and corrected sign parities}
\begin{proof}
Given a single-marker LLPO input, if its nonzero coordinate has
position \(2j\), associate \(x=2^{-j}\); if it has position
\(2j+1\), associate \(x=-2^{-j}\); with no marker use \(x=0\).
This is a computable scalar real transformation. For precision \(n\),
inspect positions through \(2n+1\). If the marker is found, output
its exact associated rational; otherwise output zero. Any later marker
has \(j\geq n+1\) and magnitude at most \(2^{-n-1}\), so every
issued rational meets the Cauchy error clause. An even marker permits
only LLPO answer one and gives a positive real, whose weak-comparison
answer is also one. An odd marker similarly permits only zero and
gives a negative real. No marker permits both in each presentation.
Copying the real relation's answer proves one strong reduction.

Conversely, on a Cauchy name \((u_k)\) of \(x\), dovetail the
strict tests \(u_k-2^{-k}>0\) and \(u_k+2^{-k}<0\).
At the first positive certificate schedule one marker at an even
position; at the first negative certificate schedule one marker at
an odd position. Until a certificate and after that one marker issue
zeros. Both tests cannot succeed on one valid name. For a nonzero
real its true sign is eventually certified, and for zero neither is.
Thus the generated stream is always a valid single-marker input and
its allowed answers are exactly the source relation's sign answers.
The oracle answer is copied independently of the original real name.
This proves the reverse strong reduction, without deciding equality
or requiring a fixed time bound near zero.
\end{proof}
\end{annexitem}

\begin{annexitem}{annex:p04-wc-weak-single}{thm:p04-wc-weak-single}{Recovering each unique output digit by a finite cover}
\begin{proof}
Let \(K,H\) be ordinary reduction operators, and fix a promised input
stream \(p\). The preprocessor supplies a negative name of a
nonempty closed \(A_p\subseteq\mathcal C\).
For every \(b\in A_p\), the computation
\(H\langle p,b\rangle\) returns exactly the Baire stream \(f(p)\),
because its output representation is identity and \(f\) is single-valued.

Fix an output coordinate \(n\). For each candidate natural value
\(v\), enumerate finite binary cylinders on which the computation
of that coordinate, with \(p\) as its other input, completes with
value \(v\). A completed finite computation reads only a finite
binary prefix, so its cylinder is a valid certificate for all
extensions. Dovetail finite-prefix simulations to enumerate their
union \(U_v\). This is uniform relative to the given input stream
\(p\), and requires no oracle besides that input.
Now test whether the enumerated cylinders of \(U_v\), together with
the negatively enumerated complement of \(A_p\), cover Cantor space.
The finite common-depth test in the preceding proof makes this a
semidecision by finite covers.

For \(v=f(p)(n)\), every \(b\in A_p\) completes the required
coordinate with that value, so \(A_p\subseteq U_v\).
Compactness supplies a finite subcover after adjoining the exterior;
the test eventually succeeds. For any different \(v\), no element
of \(A_p\) lies in \(U_v\), so the test cannot succeed, since
\(A_p\) is nonempty. Dovetailing all candidates therefore computes
the unique digit. Repeating for each coordinate gives a productive
computable realizer of \(f\) on every promised input name.
The argument does not assume that names of the same non-Baire
output object have identical digits; that would invalidate this
particular proof outside its stated output representation.
\end{proof}
\end{annexitem}

\begin{annexitem}{annex:p04-wc-no-computable-point}{thm:p04-wc-no-computable-point}{A small effective cylinder cover and discrete-answer separation}
\begin{proof}
Use the fixed numerical program enumeration from \cref{thm:enumeration}.
For program index \(e\), dovetail its computations on inputs
\(0,\ldots,e+1\). If all halt and their outputs are binary, enumerate
the cylinder of that length-\(e+2\) word, at most once for this index.
Every computable binary stream is covered, since some program computes
all its entries, including this finite prefix.
The sum of the possible cylinder weights is
\(\sum_{e=0}^{\infty}2^{-e-2}=1/2\).
No finite subcollection covers Cantor space: at common depth \(L\)
the number of covered binary words is at most
\(2^L\sum_e2^{-e-2}\leq2^{L-1}<2^L\), by counting each cylinder's
extensions and allowing overlap only to increase this upper bound.
Compactness therefore rules out coverage by the entire open enumeration.
Its complement \(A\) is nonempty, negatively computable, closed, and
has no computable point. Padding makes the enumeration a productive
negative name even at stages without a new cylinder.

If Cantor choice reduced ordinarily to natural choice, apply the
computable preprocessor to this computable negative name. It gives
a computable negative name of a nonempty natural set. Pick any allowed
number \(n\) and an oracle realizer returning a canonical computable
name for it at that input. The ordinary postprocessor now has two
computable inputs and must produce a computable element of \(A\),
contradiction. The choice of this one \(n\) need not be uniform:
every particular natural number has a computable canonical name.
For an LPO target, its preprocessed input is likewise computable,
and its one discrete answer has a computable canonical name; the same
contradiction applies. This separates an instance with no computable
solution from targets whose every discrete answer is individually
computable, without claiming that the target algorithms are computable.
\end{proof}
\end{annexitem}

\begin{annexitem}{annex:p04-wc-degree-boundaries}{thm:p04-wc-degree-boundaries}{The positive diagram and all claimed missing arrows}
\begin{proof}
First construct \(\mathsf{LPO}\leq_{\mathrm{sW}}\mathsf{C}_{\mathbb N}\).
For input \(p\), negatively name
\[
 A_p=\bigl(\{0\}\text{ if }(\forall j)\ p(j)\ne0\bigr)
       \ \cup\ \{j+1:p(j)=0\}.
\]
Exclude \(j+1\) when the finite observation \(p(j)\ne0\) is
made, and exclude zero as soon as any zero coordinate is found.
This is a computable complement enumeration, not a decision of the
universal condition. The set is always nonempty. Its only allowed
answer in the no-zero case is zero, which decodes to LPO value one;
every allowed answer in the other case is positive and decodes to
LPO value zero. Only the chosen number is needed in postprocessing.

Binary choice embeds strongly into Cantor choice by mapping a nonempty
binary set to the streams with permitted first bit and zero tail;
enumerate cylinders for an excluded first bit and all cylinders with
a nonzero later bit. Equivalently use extensiveness and
\(\mathsf{LLPO}\leq_{\mathrm{sW}}\widehat{\mathsf{LLPO}}\).
Thus \(\mathsf{LLPO}\leq_{\mathrm W} \mathsf{C}_{\mathcal C}\).

Native LPO is single-valued with a discrete answer. If it reduced
ordinarily to Cantor choice, composing its postprocessor with
canonical padding of the bit would give a single-valued Baire-output
problem \(p\mapsto(\mathsf{LPO}(p),0,0,\ldots)\) below Cantor choice.
The weak-single theorem would make this computable, contradicting
the LPO continuity obstruction. Hence
\(\mathsf{LPO}\not\leq_{\mathrm W} \mathsf{C}_{\mathcal C}\).
It also cannot reduce to LLPO, since the latter reduces to Cantor choice;
the positive LLPO-to-LPO reduction is therefore strict.
Because \(\mathsf{LPO}\leq_{\mathrm W} \mathsf{C}_{\mathbb N}\), one also has
\(\mathsf{C}_{\mathbb N}\not\leq_{\mathrm W} \mathsf{C}_{\mathcal C}\).
The no-computable-point theorem proves the other incomparability arrows
and \(\mathsf{C}_{\mathcal C}\not\leq_{\mathrm W}\mathsf{LPO}\).

Parallelizing \(\mathsf{LLPO}\leq_{\mathrm{sW}}\mathsf{LPO}\) and using
the established equivalences gives
\(\mathsf{C}_{\mathcal C}\leq_{\mathrm{sW}}\lim\).
If the reverse ordinary reduction held, then extensiveness
\(\mathsf{LPO}\leq_{\mathrm W}\widehat{\mathsf{LPO}}\equiv_{\mathrm W}\lim\)
would give the forbidden LPO-to-Cantor-choice reduction.
Thus \(\mathsf{C}_{\mathcal C}<_{\mathrm W}\lim\).
Finally \(\mathsf{C}_{\mathbb N}\leq_{\mathrm{sW}}\widehat{\mathsf{C}_{\mathbb N}}
\equiv_{\mathrm{sW}}\lim\) follows by duplication. Every claimed positive
arrow and every strictness or incomparability in the statement is
therefore supported by a specific construction or contradiction.
\end{proof}
\end{annexitem}
